\documentclass[conference]{IEEEtran}
\IEEEoverridecommandlockouts

\usepackage{cite}
\usepackage{amsmath,amssymb}
\usepackage{graphicx}
\usepackage{xcolor}
\usepackage{booktabs}
\usepackage{hyperref}
\usepackage{placeins}

% Allow more floats per page
\renewcommand{\topfraction}{0.95}
\renewcommand{\bottomfraction}{0.85}
\renewcommand{\textfraction}{0.05}
\renewcommand{\floatpagefraction}{0.80}
\setcounter{topnumber}{4}
\setcounter{bottomnumber}{4}
\setcounter{totalnumber}{8}

\graphicspath{{../outputs/results/}}

\begin{document}

\title{Asymmetric Sentiment Effects in Earnings Calls:\\
Predicting Abnormal Returns with NLP}

\author{%
\IEEEauthorblockN{Borja Valentin Albers-Sch\"{o}enberg Engeler}
\IEEEauthorblockA{\textit{School of Science and Technology}\\
IE University, Madrid, Spain}
\and
\IEEEauthorblockN{Juan Alonso-Allende Zabala}
\IEEEauthorblockA{\textit{School of Science and Technology}\\
IE University, Madrid, Spain}
\and
\IEEEauthorblockN{Javier Dom\'{i}nguez Segura}
\IEEEauthorblockA{\textit{School of Science and Technology}\\
IE University, Madrid, Spain}
\and
\IEEEauthorblockN{Alejandro Helmrich Laura}
\IEEEauthorblockA{\textit{School of Science and Technology}\\
IE University, Madrid, Spain}
\and
\IEEEauthorblockN{Luis Andr\'{e}s Infante N\'{u}\~{n}ez}
\IEEEauthorblockA{\textit{School of Science and Technology}\\
IE University, Madrid, Spain}
\and
\IEEEauthorblockN{Juliette Janne Dothee}
\IEEEauthorblockA{\textit{School of Science and Technology}\\
IE University, Madrid, Spain}
\and
\IEEEauthorblockN{Diego Oliveros Rabago}
\IEEEauthorblockA{\textit{School of Science and Technology}\\
IE University, Madrid, Spain}
\and
\IEEEauthorblockN{Lucas Karl Van Zyl}
\IEEEauthorblockA{\textit{School of Science and Technology}\\
IE University, Madrid, Spain}
}

\maketitle

% ============================================================
\begin{abstract}
Earnings calls are pre-announced corporate events that produce rich unstructured
text containing tone and sentiment information beyond what headline financial
metrics capture. This paper tests the \emph{asymmetry hypothesis}: that negative
language is a disproportionately stronger predictor of short-window Cumulative
Abnormal Returns (CAR) than positive language, consistent with investor loss aversion.
We build a corpus of 620 earnings call transcripts from 19 technology-sector
firms spanning 2007--2025 and pair each transcript with event-study CAR targets
derived from an OLS market model. Sentiment features are extracted using the
Loughran--McDonald financial lexicon, TF-IDF bag-of-words, and Word2Vec
document embeddings. A Wald test rejects the symmetry null at the 5\% level
($p = 0.042$) in the LM + controls specification, with $\hat\beta_{\text{neg}} = -8.92$
versus $\hat\beta_{\text{pos}} = +3.31$, a 2.7$\times$ asymmetry ratio. The
effect is concentrated in the Q\&A section, where management language is least
scripted. The LM lexicon achieves the best out-of-sample $R^2$ (+0.057),
outperforming TF-IDF and Word2Vec despite its simplicity. These results confirm
that domain-adapted sentiment features carry genuine predictive information over
market baselines.
\end{abstract}

\begin{IEEEkeywords}
Natural Language Processing, earnings calls, sentiment analysis, event study,
Loughran--McDonald lexicon, asymmetry hypothesis, abnormal returns, TF-IDF, Word2Vec
\end{IEEEkeywords}

% ============================================================
\section{Introduction}
\label{sec:intro}
% ============================================================

Every quarter, publicly traded companies host earnings calls in which executives
present financial results and respond to questions from equity analysts.
These events are pre-announced and occur on a fixed schedule, which makes them
a natural controlled setting for studying the link between corporate language
and short-window market reactions. The text they produce (prepared remarks,
forward-looking guidance, and unscripted Q\&A exchanges) contains information
that goes well beyond the headline numbers.

NLP has made it practical to extract structured sentiment signals from large
transcript corpora at low cost, using methods ranging from domain-adapted
lexicons to neural embeddings. What is less clear is whether those signals are
genuinely informative about subsequent returns, and in particular whether
\emph{negative} language carries disproportionately more weight than equally
salient positive language. Behavioral finance research on loss aversion and
negativity bias~\cite{ref_suzuki2020} suggests that markets should react more
strongly to bad news, but this has not been formally tested on earnings call
transcripts with the statistical rigor of a Wald test.

This paper addresses that gap. Our core hypothesis is:
\begin{equation}
H_1: |\beta_{\text{neg}}| > |\beta_{\text{pos}}|
\end{equation}
where $\beta_{\text{neg}}$ and $\beta_{\text{pos}}$ are OLS coefficients on
negative and positive sentiment rates when predicting CAR following earnings calls.
We test it formally via a Wald test of the joint symmetry constraint
$H_0: \beta_{\text{neg}} + \beta_{\text{pos}} = 0$.

What is unique about our approach is the combination of three elements: (i) a
section-level sentiment decomposition that separates the scripted presentation
from the unscripted Q\&A, (ii) a rigorous out-of-sample $R^2$ evaluation against
a null forecast, and (iii) a formal statistical test of asymmetry rather than
informal comparison of coefficients. Our dataset spans 19 technology-sector
firms over 2007--2025 (620 events), large enough to detect the effect while
keeping the sample homogeneous.

Our main finding is that the asymmetry hypothesis is supported at the 5\%
significance level (Wald $p = 0.042$) with a 2.7$\times$ ratio
($\hat\beta_{\text{neg}} = -8.92$ vs.\ $\hat\beta_{\text{pos}} = +3.31$).
The effect is driven entirely by the Q\&A section of each call; the prepared
presentation carries essentially no negative signal. The simple two-feature
LM lexicon achieves the best OOS $R^2$ of $+0.057$, outperforming TF-IDF
and Word2Vec.

The remainder of the paper is organized as follows. Section~\ref{sec:litreview}
reviews related work on financial text mining and return prediction.
Section~\ref{sec:methodology} describes the data, event study design, and NLP
pipeline. Section~\ref{sec:results} presents and interprets all results.
Section~\ref{sec:deployment} covers the interactive Streamlit deployment.
Section~\ref{sec:conclusion} concludes and Section~\ref{sec:futurework}
identifies directions for future work.\footnote{Detailed sub-reports for each
pipeline stage (data collection, preprocessing, feature extraction, modelling,
evaluation, deployment, and code documentation) are available at
\url{https://github.com/Tee-Time-Software-Solutions/NASDAQ-NLP/tree/main/docs/report}.}

% ============================================================
\section{Related Literature}
\label{sec:litreview}
% ============================================================

The question of whether financial text predicts returns has been studied from
several angles. We draw on three closely related works.

\subsection{Opinion vs.\ fact decomposition for forecasting}

Suzuki et al.~\cite{ref_suzuki2020} build a pipeline on analyst reports that
first separates \emph{opinion} sentences from factual ones, then uses those
signals to forecast net income estimate changes and stock price movements.
Their key finding is that combining text signals with numeric trend features
improves fundamental forecasting meaningfully, but that predicting price
movements directly from the same text is much harder, on some targets the
model approaches random-baseline performance.

This work established realistic expectations for our project: if sentiment-to-return
forecasting is hard even with carefully engineered NLP, modest OOS $R^2$ values
are not a failure. It also motivated our decision to split transcripts into
presentation and Q\&A sections, treating earnings text as heterogeneous rather
than a single undifferentiated document.

\subsection{Transformer-based sentiment and return predictability}

Schlaubitz~\cite{ref_schlaubitz2021} applies a fine-tuned DistilBERT model to
earnings reports for Swiss SMEs and studies the link between extracted sentiment
and subsequent returns. The relationship is noisy and heterogeneous across
firms; adding historical market features helps in nonlinear settings, but the
Transformer approach does not consistently outperform simpler lexicon-based
methods on small, domain-specific datasets.

This finding shaped our evaluation framing: we judge sentiment features by their
\emph{incremental predictive value} over market-only baselines, not by raw
accuracy. We also originally planned to include FinBERT as a feature source,
but Transformer inference over 620 long transcripts exhausted local memory,
making it infeasible. Schlaubitz's results suggest this is not a critical omission.

\subsection{NLP pipelines and domain challenges in financial sentiment}

Omoseebi et al.~\cite{ref_omoseebi2023} survey NLP pipelines for financial
sentiment analysis, covering preprocessing choices, TF-IDF vs.\ embedding
representations, and model families. The central challenge they identify is
specialized vocabulary: words like ``loss,'' ``liability,'' and ``default'' are
clearly negative in financial contexts but would be scored neutrally by
general-purpose sentiment tools.

This motivated our choice of the Loughran--McDonald (LM) dictionary over VADER
or AFINN. The LM lexicon was built specifically for financial document analysis
and has been validated on SEC filings and earnings materials. Our pipeline
structure (acquisition $\rightarrow$ preprocessing $\rightarrow$ features
$\rightarrow$ model $\rightarrow$ OOS evaluation) follows the canonical design
they describe.

% ============================================================
\section{Experimental Methodology}
\label{sec:methodology}
% ============================================================

\subsection{Part 1: Algorithms and Representations}

\subsubsection{Event Study Market Model}

The event study framework~\cite{ref_suzuki2020} uses OLS to estimate each
stock's expected return from a pre-event estimation window and then measures
the surprise (abnormal return) around the event. For a stock $i$ and market
index $m$:
\begin{equation}
R_{i,t} = \alpha_i + \beta_i \, R_{m,t} + \varepsilon_{i,t}, \quad t \in [-120,-20]
\end{equation}
The abnormal return on day $t$ is the out-of-sample residual:
\begin{equation}
AR_{i,t} = R_{i,t} - \bigl(\hat\alpha_i + \hat\beta_i \, R_{m,t}\bigr)
\end{equation}
Cumulative abnormal returns are:
\begin{align}
\text{CAR}_{[0,1]} &= AR_0 + AR_1 \\
\text{CAR}_{[0,3]} &= AR_0 + AR_1 + AR_2 + AR_3
\end{align}
CAR$_{[0,3]}$ is our primary target because it captures both the immediate
reaction and any post-announcement drift.

\subsubsection{Loughran--McDonald Lexicon}

The LM Master Dictionary~\cite{ref_omoseebi2023} contains 2,345 \emph{Negative}
and 347 \emph{Positive} financial-domain words. For each transcript section $s$:
\begin{equation}
\text{NegRate}_s = \frac{N_{\text{neg},s}}{|W_s|}, \quad
\text{PosRate}_s = \frac{N_{\text{pos},s}}{|W_s|}
\end{equation}
Normalizing by word count $|W_s|$ controls for transcript length.
We compute these for three scopes (full transcript, presentation, Q\&A)
giving six features per event.

\subsubsection{TF-IDF}

Term Frequency--Inverse Document Frequency represents each document as a
weighted sparse vector. We use a vocabulary of 500 terms, minimum document
frequency 3, unigrams only. The vocabulary is fitted on the training set
exclusively to avoid look-ahead bias.

\subsubsection{Word2Vec}

A skip-gram Word2Vec model~\cite{ref_suzuki2020} (100 dimensions, window 5,
10 epochs) is trained on training-set transcripts. Each document is represented
as the mean of its in-vocabulary token embeddings, yielding a dense 100-dimensional
vector. Tokens absent from the vocabulary are skipped during pooling.

\subsubsection{Regression and Classification Models}

For regression (predicting CAR$_{[0,3]}$) we evaluate OLS with HC3
standard errors, Ridge ($\alpha = 1$), Random Forest (100 trees), and MLP.
OOS $R^2$ follows Campbell and Thompson (2008):
\begin{equation}
\text{OOS }R^2 = 1 - \frac{\sum_t (y_t - \hat{y}_t)^2}{\sum_t (y_t - \bar{y}_{\text{train}})^2}
\end{equation}
Positive OOS $R^2$ means the model beats a naive historical-mean forecast.

The asymmetry hypothesis is tested via a Wald test on the OLS coefficient vector:
\begin{equation}
H_0: \beta_{\text{neg}} + \beta_{\text{pos}} = 0 \quad \text{(symmetric effects)}
\end{equation}

For classification (predicting whether CAR$_{[0,3]} > 0$) we evaluate
Naive Bayes, Logistic Regression, Decision Tree, Random Forest, and MLP,
reporting accuracy, F1, and AUC.

\subsection{Part 2: Implementation}

\subsubsection{Data and Corpus}

Our corpus combines two sources (Table~\ref{tab:dataset}).
The primary set is 188 Thomson Reuters StreetEvents transcripts for 10
NASDAQ-listed firms (2016--2020). A supplementary Kaggle corpus
(\texttt{cleaned\_ECTs}) adds 9 technology-adjacent firms including IBM,
Oracle, Salesforce, and Accenture, extending coverage to 2007--2025.
The sample is restricted to the technology sector to minimize cross-sector
heterogeneity in the sentiment-return relationship; 620 events remain after
applying this filter. Daily return data is downloaded from Yahoo Finance via
\texttt{yfinance}. NASDAQ-listed stocks use the Composite (\texttt{\^{}IXIC})
as their benchmark; NYSE-listed stocks use the S\&P 500 (\texttt{\^{}GSPC}).

\begin{table}[htbp]
\caption{Dataset Summary}
\label{tab:dataset}
\centering
\begin{tabular}{ll}
\toprule
Attribute & Value \\
\midrule
Total events & 620 \\
Unique tickers & 19 \\
Date range & 2007--2025 \\
Training window & 2016--2018 ($n \approx 155$) \\
Test window & 2019--2020 ($n \approx 187$) \\
Market benchmark (NASDAQ) & \^{}IXIC \\
Market benchmark (NYSE) & \^{}GSPC \\
\bottomrule
\end{tabular}
\end{table}

\subsubsection{Event Trading Day Assignment}

A critical implementation detail is mapping each call to the correct trading day.
Calls ending before 4:00 PM ET are assigned to the same calendar day; calls
at or after market close are pushed to the next business day. Getting this wrong
would shift the event window by one day and corrupt the CAR measure entirely.

\subsubsection{Text Preprocessing}

Each transcript is processed through four stages:
\begin{enumerate}
  \item Remove operator instructions, legal disclaimers, and safe-harbor boilerplate.
  \item Split into \textbf{presentation} and \textbf{Q\&A} sections using
        participant turn markers.
  \item Lowercase and standardize punctuation; retain stopwords and unstemmed
        forms for unambiguous LM matching.
  \item Tokenize with \texttt{[A-Za-z']+}, consistent with LM dictionary boundaries.
\end{enumerate}

\subsubsection{Evaluation Protocol}

All models use a strict time-based split: 2016--2018 training, 2019--2020 test.
No test-period data enters feature vocabulary construction or hyperparameter
selection. An expanding-window cross-validation (train on 2016$\rightarrow$Y,
test on Y$+$1, for Y $\in$ \{2017, 2018, 2019\}) provides additional robustness.

\subsubsection{Deployment}

The pipeline is packaged as an installable Python library (\texttt{nasdaq-nlp},
managed with \texttt{uv}) and deployed as a Streamlit dashboard accessible
at \url{https://nasdaq-nlp.streamlit.app/}.
The dashboard serves pre-computed CSVs from four tabs mirroring the four
analysis notebooks, with Altair charts and a global sidebar for ticker and
CAR window selection.\footnote{Full source code and notebooks:
\url{https://github.com/Tee-Time-Software-Solutions/NASDAQ-NLP}.}

% ============================================================
\section{Results and Discussion}
\label{sec:results}
% ============================================================

\subsection{Asymmetry Test}

Table~\ref{tab:asymmetry} reports Wald test results across OLS specifications.
The plain LM model ($\hat\beta_{\text{neg}} = -3.51$, Wald $p = 0.485$)
does not reject $H_0$. Adding pre-event volatility and year fixed effects
changes the picture substantially: the negative coefficient grows to $-8.92$
and $H_0$ is rejected at the 5\% level ($p = 0.042$). The asymmetric signal
is not in the raw sentiment rates, it is in the residual after controlling
for the macro environment.

The Q\&A section drives the effect. The presentation section alone produces
$\hat\beta_{\text{neg}} \approx -0.07$, consistent with sanitized executive
language in prepared remarks. The Q\&A coefficient reaches $-10.3$ in the
Pres + Q\&A model, confirming that unscripted exchanges contain the
decision-relevant negative information.

\begin{table}[htbp]
\caption{Asymmetry Test Results (CAR$_{[0,3]}$)}
\label{tab:asymmetry}
\centering
\begin{tabular}{lrrrr}
\toprule
Model & $\hat\beta_{\text{neg}}$ & $\hat\beta_{\text{pos}}$ & Ratio & Wald $p$ \\
\midrule
LM Lexicon               & $-3.51$ & $+2.39$ & $1.5\times$ & $0.485$ \\
LM + Controls [OLS]      & $-8.92$ & $+3.31$ & $2.7\times$ & $\mathbf{0.042}$ \\
LM Pres only$^\dagger$   & $-0.07$ & $+0.64$ & ---         & $0.594$ \\
LM Q\&A only$^\dagger$   & $-4.69$ & $+3.66$ & $1.3\times$ & $0.386$ \\
LM Pres + Q\&A$^\ddagger$ & $-10.3$ & $+3.08$ & $3.3\times$ & $0.218$ \\
\bottomrule
\end{tabular}
{\footnotesize
$^\dagger$Section-specific rates used instead of full-transcript rates.\\
$^\ddagger$Q\&A coefficients shown; presentation yields
$\hat\beta_{\text{neg,pres}} = +1.82$ (no negative signal in scripted remarks).}
\end{table}

\begin{figure*}[!t]
\centering
\includegraphics[width=\textwidth]{coefficient_plot.png}
\caption{OLS sentiment coefficient estimates ($\hat\beta_{\text{neg}}$, $\hat\beta_{\text{pos}}$)
for each LM specification on CAR$_{[0,3]}$. Red bars indicate individual significance at
$p < 0.10$. The LM Pres+Q\&A panel illustrates the section-level contrast:
scripted presentation language ($\hat\beta_{\text{neg,pres}} = +1.82$) carries no
negative signal, while Q\&A ($\hat\beta_{\text{neg,qa}} = -10.3$) drives the
asymmetric result.}
\label{fig:coef}
\end{figure*}

\FloatBarrier
\subsection{Out-of-Sample Predictive Performance}

Table~\ref{tab:oos} reports regression OOS $R^2$ on the 2019--2020 test window.
The simple two-feature LM lexicon achieves the best OOS $R^2$ of $+0.057$,
outperforming all richer alternatives. Word2Vec embeddings reach $+0.050$
and TF-IDF with Random Forest $+0.045$.

Models with year dummies and Ridge regularization systematically underperform:
LM + Controls [Ridge] gives $-0.044$ versus $-0.086$ for OLS, showing that
regularization partially mitigates but cannot eliminate overfitting with
$\sim$155 training observations and many fixed-effect parameters.
TF-IDF + OLS ($-0.469$) is the most extreme case of this overfitting pattern.

\begin{table}[htbp]
\caption{Out-of-Sample $R^2$: Regression on CAR$_{[0,3]}$}
\label{tab:oos}
\centering
\begin{tabular}{lr}
\toprule
Model & OOS $R^2$ \\
\midrule
Null (historical mean)          & $0.000$ \\
\textbf{LM Lexicon [OLS]}       & $\mathbf{+0.057}$ \\
Word2Vec [Ridge]                 & $+0.050$ \\
TF-IDF [RF]                     & $+0.045$ \\
TF-IDF [Ridge]                  & $+0.043$ \\
LM + Controls [Ridge]           & $-0.044$ \\
LM + Controls [OLS]             & $-0.086$ \\
TF-IDF [OLS]                    & $-0.469$ \\
\bottomrule
\end{tabular}
\end{table}

\begin{figure*}[!t]
\centering
\includegraphics[width=\textwidth]{oos_vs_clf_plot.png}
\caption{Left: OOS $R^2$ across regression models (red bars below null baseline).
Right: Test-set accuracy and F1 for direction classifiers (dashed = 50\% random baseline).
The LM lexicon leads regression despite its simplicity; TF-IDF models lead classification.}
\label{fig:oos_clf}
\end{figure*}

\FloatBarrier
\subsection{Robustness Across Event Windows}

Table~\ref{tab:rob} replicates the LM experiments with CAR$_{[0,1]}$ as the
target. The sign of $\hat\beta_{\text{neg}}$ remains negative in both windows,
confirming directional consistency. The Wald $p$-value is higher for the shorter
window ($p = 0.779$ vs.\ $0.485$), suggesting the asymmetric signal accumulates
over multiple post-event days rather than being priced entirely on day zero.

\begin{table}[htbp]
\caption{Robustness: LM Lexicon Across Event Windows}
\label{tab:rob}
\centering
\begin{tabular}{llrr}
\toprule
Model & Window & Wald $p$ & OOS $R^2$ \\
\midrule
LM Lexicon & CAR$_{[0,3]}$ & $0.485$ & $+0.057$ \\
LM Lexicon & CAR$_{[0,1]}$ & $0.779$ & $+0.052$ \\
LM Q\&A    & CAR$_{[0,3]}$ & $0.386$ & --- \\
LM Q\&A    & CAR$_{[0,1]}$ & $0.692$ & --- \\
\bottomrule
\end{tabular}
\end{table}

\FloatBarrier
\subsection{Classification Results}

Table~\ref{tab:clf} reports direction-prediction metrics. The best AUC is
$0.600$ for TF-IDF + LM + Controls [RF]; the best accuracy is $59.6\%$,
achieved by both TF-IDF [NaiveBayes] and TF-IDF + LM + Controls [RF].
All TF-IDF-based models exceed the 50\% random baseline in accuracy.

LM-only classifiers show lower accuracy but higher F1, reflecting a class
imbalance in the test window: positive CAR events dominated the 2019--2020
bull market, so a model that predicts positive more aggressively gets higher F1
at the cost of accuracy.

\begin{table}[htbp]
\caption{Classification Results (CAR$_{[0,3]} > 0$, Test 2019--2020)}
\label{tab:clf}
\centering
\begin{tabular}{lrrr}
\toprule
Model & Accuracy & F1 & AUC \\
\midrule
TF-IDF + LM + Controls [RF] & $59.6\%$ & $0.553$ & $\mathbf{0.600}$ \\
TF-IDF [NaiveBayes]         & $\mathbf{59.6\%}$ & $0.632$ & $0.588$ \\
TF-IDF [RF]                  & $51.9\%$ & $0.468$ & $0.563$ \\
TF-IDF [LogReg]              & $51.9\%$ & $0.590$ & $0.578$ \\
LM Lexicon [LogReg]          & $45.5\%$ & $0.625$ & $0.508$ \\
\bottomrule
\end{tabular}
\end{table}

\begin{figure*}[!t]
\centering
\includegraphics[width=\textwidth]{benchmark_plots.png}
\caption{Six-panel benchmark. Top: OOS $R^2$ and train/test $R^2$ overfitting check.
Middle: MAE and Wald $p$-values per OLS specification.
Bottom: Classifier accuracy/F1 and train/test accuracy overfitting check.
All panels evaluated on CAR$_{[0,3]}$, test window 2019--2020.}
\label{fig:benchmark}
\end{figure*}

\FloatBarrier
\subsection{Discussion}

The results support the asymmetry hypothesis and tell a consistent story.
Negative sentiment in earnings calls is not a noisy proxy, it carries
genuine predictive information, but only after controlling for background market
risk and year effects. The mechanism is intuitive: management reads sanitized,
legally reviewed prepared remarks, so the presentation section's negative rate
carries almost no signal ($\hat\beta_{\text{neg,pres}} \approx 0$). The Q\&A
exchange is different, analysts ask pointed questions and management responds
without a script, which is where real tone information surfaces ($\hat\beta_{\text{neg,qa}} = -10.3$).

The simplicity premium in OOS $R^2$ is expected in this setting.
With $\sim$155 training events, a 500-feature TF-IDF matrix is severely
underdetermined and Ridge regularization can only partially compensate.
Domain-adapted features that are noise-resistant by design outperform
representation learning in this low-data regime, which is consistent with
Schlaubitz's finding that DistilBERT does not consistently outperform simpler
approaches on small domain-specific datasets~\cite{ref_schlaubitz2021}.

The 60\% AUC ceiling on classification is also not surprising. Market reactions
are partially driven by factors orthogonal to language (algorithmic order flow,
options market hedging, and analyst surprise relative to consensus forecasts),
none of which are captured in transcript text alone. A ceiling around 60\%
accuracy is expected and is not a failure of the NLP approach.

One potential objection is that the asymmetry result ($p = 0.042$) could be
driven by a small number of events with extreme negative language coinciding
with large market moves. The robustness of $\hat\beta_{\text{neg}} < 0$ across
both event windows and the Q\&A-specific decomposition argue against this,
but a larger and more diverse corpus would further strengthen the claim.

% ============================================================
\section{Deployment}
\label{sec:deployment}
% ============================================================

The pipeline is deployed as an interactive Streamlit dashboard
(\url{https://nasdaq-nlp.streamlit.app/})
organized into four tabs mirroring the four analysis notebooks:
Data Pipeline, Feature Extraction, Modelling, and Results.
A fifth ``Try It Yourself'' tab accepts free text input and returns
a live LM sentiment score and predicted CAR direction.
Pre-computed CSVs are committed to \texttt{frontend/data/} so the app
loads without re-executing the pipeline, achieving sub-second response times.
All charts use Altair for interactive tooltips; a global sidebar propagates
ticker and CAR window selection to all tabs simultaneously.

% ============================================================
\section{Conclusion}
\label{sec:conclusion}
% ============================================================

We have shown that negative sentiment in earnings call transcripts is a
statistically stronger predictor of short-window abnormal returns than positive
sentiment, supporting the asymmetry hypothesis at the 5\% level
($p = 0.042$, asymmetry ratio $2.7\times$). The effect is concentrated in the
Q\&A section and disappears in the scripted presentation, providing a clear
mechanism: unscripted exchanges with analysts are where honest management tone
is revealed.

The simple Loughran--McDonald lexicon achieves the best out-of-sample predictive
performance (OOS $R^2 = +0.057$), outperforming TF-IDF and Word2Vec.
This is a practically important finding: domain-adapted, low-dimensional features
beat representation learning when training data is limited to $\sim$155 firm-level
events. The classifiers reach 60\% AUC, modestly above chance, which is
consistent with the view that a large fraction of post-earnings return variance is
driven by factors outside the transcript text.

% ============================================================
\section{Future Work}
\label{sec:futurework}
% ============================================================

Three extensions stand out as the most promising next steps.

First, incorporating \textbf{analyst consensus data} (earnings beat/miss vs.\
expectations) as a control variable would isolate the pure tone signal from the
fundamental surprise component, which is the dominant driver of post-earnings
CAR and is currently absent from our specification.

Second, \textbf{expanding to a multi-sector corpus} would test whether the
asymmetry effect is specific to technology firms or holds more broadly across
industries with different earnings communication styles.

Third, integrating \textbf{FinBERT}~\cite{ref_omoseebi2023}, a BERT model
fine-tuned on financial text, as an additional feature source is a natural
extension. We attempted this during development but Transformer inference over
620 transcripts of up to 15,000 tokens each exhausted local memory resources.
Dedicated GPU compute would make this feasible and would likely improve
classification performance in particular.

% ============================================================
\section*{Acknowledgements}
% ============================================================

The authors thank IE University for providing access to course materials
and computational resources. Transcript data was sourced from Thomson Reuters
StreetEvents and the Kaggle \texttt{cleaned\_ECTs} corpus; market data was
obtained from Yahoo Finance via the \texttt{yfinance} library.

% ============================================================
\begin{thebibliography}{00}

\bibitem{ref_suzuki2020}
M. Suzuki, H. Sakaji, K. Izumi, H. Matsushima, and Y. Ishikawa,
``Forecasting Net Income Estimate and Stock Price Using Text Mining from
Economic Reports,''
\emph{Information}, vol.~11, no.~292, 2020.

\bibitem{ref_schlaubitz2021}
A. Schlaubitz,
\emph{Natural Language Processing in Finance: Analysis of Sentiment and
Complexity of News and Earnings Reports of Swiss SMEs and their Relevance
for Stock Returns}.
Master's Thesis, Zurich University of Applied Sciences (ZHAW), 2021.

\bibitem{ref_omoseebi2023}
A. Omoseebi, J. Owen, and Q. Loveth,
``Natural language processing (NLP) for sentiment analysis in
financial markets,'' 2023.

\end{thebibliography}

\end{document}
